\section{Introduction}

Phishing email is still one of the most common ways into an organization. The
attacker does not break a technical control. They write a message that looks
like it came from a bank, a colleague or a supplier, and the reader hands over
the password. Generative AI made writing such mail cheap: Heiding et
al.~\cite{heiding} showed that phishing written by GPT-4 already gets click
rates close to phishing written by human experts.

Detection results in the literature look excellent. ChatSpamDetector reports
99.70\% accuracy with GPT-4~\cite{koide}, a fine-tuned RoBERTa reports
98.45\%~\cite{uddin}, and a five agent LLM system reports 97.89\%~\cite{multiphish}.
When we reproduced the setting with TF-IDF and logistic regression we also got
97 to 99\%. A method older than deep learning matching GPT-4 is a warning sign:
it suggests the benchmark, not the model, is doing the work.

This paper asks what these detectors are worth when the test mail comes from a
corpus the model has never seen, and when the emails that the corpora share
with each other are removed first. Three things make the question concrete.
First, the popular corpora are not independent. Aggregated sets such as the
Kaggle phishing email set are built from older collections, so the same email
can sit on both sides of a train and test split even when the two sides come
from datasets with different names. Second, the cost of running a detector is
never reported, although for a university or a small company it decides what
can be deployed at all. Third, a single F1 score hides the thing an operator
cares about most, which is how much legitimate mail gets flagged.

Our contributions are:

\begin{enumerate}
\item A measurement of email overlap between nine public phishing and spam
sources at three levels of strictness, with every near-duplicate pair verified
by exact Jaccard similarity and the matcher itself validated against brute
force. Exact matching, the check most people would run, finds almost none of
the overlap.
\item A decontaminated leave-one-corpus-out benchmark, with a control that
removes the same number of training emails at random. The control separates
the effect of leaked emails from the effect of simply having less data, which
no earlier work on this question reports.
\item A measurement of what the corpora actually contain. An LLM annotator with
a second annotator for agreement finds that only 15 to 25\% of the positive
emails in these corpora are phishing; the rest is bulk spam.
\item A comparison on identical emails of classical models, a fine-tuned
DistilBERT, six small open LLMs, three published phishing detectors and two
current commercial models, reporting the false alarm rate and the measured
dollar cost next to accuracy.
\item An analysis of AI-written phishing that finds a give-away token in the
E-PhishLLM corpus and then shows, with a paired test on sanitised copies of the
same emails, that it is not what the detectors respond to, together with
results on the Italian and German parts of the corpus.
\item An ablation that changes only the source of the few-shot examples, which
shows that examples drawn from old corpora, and not few-shot prompting itself,
are what lower performance on AI-written phishing.
\item Operating-point analysis: how the ranking changes at a realistic amount
of phishing, and what two-stage detectors built from the same models achieve
at a fraction of the cost.
\end{enumerate}

The whole study ran on one 8\,GB laptop for under two US dollars of API fees.
Code, raw model answers, results and a notebook are public.\footnote{\url{https://github.com/Jahid15/phishing-llm-cross-corpus}}

\section{Related Work}

\textit{LLM detectors for phishing email.} Koide et al.~\cite{koide} prompt
GPT-4 with a chain of thought and parse the verdict. Heiding et
al.~\cite{heiding} use LLMs as judges of suspicious mail. Uddin et
al.~\cite{uddin} fine-tune RoBERTa and add word level explanations. Lin et
al.~\cite{lin} study small LLMs of about three billion parameters, the same
size range we use, and improve them with prompt engineering, LoRA fine-tuning
on explanations from a paid teacher model, and ensembling; they list the
missing cost analysis as a limitation. MultiPhishGuard~\cite{multiphish}
combines five agents and reports a 2.73\% false positive rate.
Phishsense-1B~\cite{phishsense} is a phishing-specific 1B model that reports
97.5\% on its own data and 70\% on a separate real-world set. Kuikel et
al.~\cite{kuikel} find that the most accurate model is not the one with the
most faithful explanations. A 2026 review~\cite{survey} covers the field from
both the attack and the defence side.

\textit{Cross-corpus evaluation.} Two recent works already test across
datasets, and we position ourselves against them rather than around them.
E-PhishGen~\cite{ephishgen} reassesses eight legacy corpora with a
leave-one-dataset-out protocol over classical models, DistilBERT and several
LLMs, and releases E-PhishLLM, a corpus of emails written by GPT-4o and
GPT-4o-mini, which we use as our AI-phishing test. Bhuiyan and
Bhuiyan~\cite{bhuiyan} study six corpora with leave-one-corpus-out training,
prevalence shift and artifact learning for classical models. Fomin~\cite{lodo}
reports the same kind of gap between pooled cross-validation and
leave-one-dataset-out for malicious prompt classifiers, and Gutierrez et
al.~\cite{gutierrez} show that much of the transfer gap to LLM-written
phishing can be recovered by recalibrating the decision threshold. Neither
E-PhishGen nor Bhuiyan and Bhuiyan measure or remove the overlap between the
corpora they train and test on, and none of these works report what the
detectors cost to run. Those two gaps, together with false alarm rates on
legitimate mail, are what this paper adds.

\textit{Leakage and deduplication.} Arp et al.~\cite{arp} list data snooping and
sampling bias among the most common pitfalls of machine learning in security.
Lee et al.~\cite{leededup} show that near-duplicates are pervasive in standard
text corpora and that removing them is needed for honest evaluation; our
matcher follows their MinHash approach~\cite{broder}. Liu et
al.~\cite{malwareleak} find near-identical samples leaking between train and
test splits in Android malware detection, the closest analogue in security to
what we find for email. For LLMs evaluated zero-shot there is the further
worry that public corpora were part of pre-training, which a recent survey
treats in detail~\cite{contamsurvey}.

\textit{Operating points.} Al Barwani et al.~\cite{ragfp} motivate their work by
the false positive burden that standalone LLM classifiers create. \c{S}ahin and
Mert~\cite{triage} propose conformal deferral for security alert triage, which
is the natural next step for the operating points we report.

\section{Data}

We use nine sources (Table~\ref{tab:data}). Six contain both classes and take
turns as the held-out corpus: SpamAssassin, CEAS-08, TREC-07, Ling-Spam and
Enron from the curated collection of Champa et al.~\cite{champa}, and the
Kaggle phishing email set used by~\cite{uddin}. Three are test only. Nazario
contains real credential phishing and Nigerian contains advance fee fraud,
both with no legitimate class. E-PhishLLM~\cite{ephishgen} contains phishing
and legitimate emails written by GPT-4o and GPT-4o-mini; we use its English
part as a test set and its Italian and German parts as a second test of
language transfer.

For every source the model input is the subject and the body joined by a new
line. Label 1 means phishing or spam and label 0 means legitimate. We drop
bodies shorter than 20 characters and exact duplicates inside a corpus.
Corpora larger than 10{,}000 emails are reduced to a stratified sample of
10{,}000 (seed 42) so that the three largest ones do not dominate training.
From each source we also draw a fixed evaluation subset of 300 emails,
balanced where both classes exist (150 for the one-class sets). Every model,
local or paid, is scored on exactly these emails.

Two properties of this data matter for reading our results. The positive class
is not the same across corpora: Section~\ref{sec:labels} measures what it
actually contains. And the Kaggle set has no subject field, so its text is
structurally unlike the other corpora, which is a difference that has nothing
to do with phishing.
